% 观测生成的信息。Omega=[0,1)，P为长度。X=floor(4 omega)取0,1,2,3；Y=g(X)=floor(X/2)=floor(2 omega)。
% sigma(X)由4个长1/4的区间生成，共2^4=16个事件；sigma(Y)只有空集、[0,1/2)、[1/2,1)、Omega四个事件。
% 事件{X=1}=[1/4,1/2)属于sigma(X)但不是Y的原子之并，故不属于sigma(Y)。
\begin{tikzpicture}[line width=.6pt,font=\fontsize{8}{10}\selectfont,>=Stealth]
\def\L{80mm}
% row X
\foreach \k/\fc in {0/FigureBlueFill,1/FigureGreenFill,2/FigureBlueFill,3/FigureGreenFill}{
  \draw[FigureStroke,fill=\fc] ({\k*\L/4},12mm) rectangle ({(\k+1)*\L/4},18mm);
  \node[FigureInk] at ({(\k+.5)*\L/4},15mm) {$X=\k$};}
\draw[FigureRed,line width=1.1pt] ({\L/4},12mm) rectangle ({\L/2},18mm);
\node[FigureRed,anchor=south] at ({3*\L/8},18.3mm) {$\{X=1\}=[1/4,1/2)$};
\node[FigureInk,anchor=east] at (-2mm,15mm) {$\sigma(X)$};
% row Y
\foreach \k/\fc in {0/FigureYellowFill,1/FigureBlueFill}{
  \draw[FigureStroke,fill=\fc] ({\k*\L/2},2mm) rectangle ({(\k+1)*\L/2},8mm);
  \node[FigureInk] at ({(\k+.5)*\L/2},5mm) {$Y=\k$};}
\node[FigureInk,anchor=east] at (-2mm,5mm) {$\sigma(Y)$};
% dashed refinement guides
\foreach \x in {0,0.25,0.5,0.75,1}{\draw[FigureMuted,dashed,line width=.4pt] ({\x*\L},8mm) -- ({\x*\L},12mm);}
% axis
\draw[FigureStroke] (0,-1mm) -- (\L,-1mm);
\foreach \x/\t in {0/0,0.25/{1/4},0.5/{1/2},0.75/{3/4},1/1}{
  \draw[FigureStroke] ({\x*\L},-1mm) -- ({\x*\L},-2.2mm);
  \node[FigureMuted,font=\fontsize{7.5}{9}\selectfont,anchor=north] at ({\x*\L},-2.4mm) {$\t$};}
\node[FigureInk,anchor=west] at (\L+2mm,-1mm) {$\omega$};
% event counts
\node[FigureInk,anchor=west] at (\L+6mm,15mm) {4个原子，$2^4=16$个事件};
\node[FigureInk,anchor=west] at (\L+6mm,5mm) {2个原子，$2^2=4$个事件};
\end{tikzpicture}
